% method: cong-thuc-kirchhoff
\begin{dieukienbox}
Dùng khi: phương trình sóng \textbf{thuần nhất} ($u_{tt}=c^2\Delta u$, không nguồn) trên $\mathbb{R}^3$.
Chỉ đúng cho $n=3$ chiều không gian. Vận tốc lan truyền là $c$.
\end{dieukienbox}

\begin{nhobox}
\textbf{Dữ liệu ban đầu:} $g(x)=u(x,0)$ (vị trí lúc đầu), $\;h(x)=u_t(x,0)$ (vận tốc lúc đầu).

\textbf{Công thức Kirchhoff} ($n=3$, vận tốc $c$) --- tích phân trên mặt cầu \textbf{bán kính $ct$}:
\[
  u(x,t)=\partial_t\!\left(\frac{1}{4\pi c^2 t}\int_{\partial B_{ct}(x)} g\,dS\right)
  +\frac{1}{4\pi c^2 t}\int_{\partial B_{ct}(x)} h\,dS .
\]
Mặt cầu $\partial B_{ct}(x)$ tâm $x$ bán kính $ct$ có diện tích $4\pi(ct)^2=4\pi c^2t^2$, nên
$\dfrac{1}{4\pi c^2t^2}\int_{\partial B_{ct}(x)}\!\cdots\,dS$ chính là \textbf{giá trị trung bình} trên mặt cầu đó.

Trường hợp $g=0$ (đứng yên lúc đầu) $\Rightarrow$ chỉ còn:
\[
  u(x,t)=\frac{1}{4\pi c^2 t}\int_{\partial B_{ct}(x)} h\,dS .
\]
\end{nhobox}

\begin{luuybox}
Đừng quên $c$: bán kính mặt cầu là $ct$ (\emph{không} phải $t$), hệ số là $\dfrac{1}{4\pi c^2t}$. Khi $c=1$ mới rút về $\dfrac{1}{4\pi t}$, mặt cầu bán kính $t$. (Đề thật hay gặp $c^2=4,9$ tức $c=2,3$.)
\end{luuybox}

\begin{meobox}
Nghiệm tại $(x,t)$ \emph{chỉ} phụ thuộc dữ liệu nằm \emph{đúng trên mặt cầu} $|y-x|=ct$ (không phải bên trong).
Vậy nếu mặt cầu này không chạm vùng mà $g,h$ khác $0$ thì $u(x,t)=0$.
\end{meobox}
